\documentclass[UTF8,a4paper,10pt]{ctexart}
\usepackage{ipara-handbook}
\usepackage{booktabs,tabularx}
\begin{document}
\section*{H-B03\quad 微分与累积：基本定理及正则性边界}
\textbf{边界：}本桥连接局部变化率与累积量，导数和积分的独立机制分别回到 Topic03、Topic05。
\begin{tcolorbox}[title=三类对象必须分开]
\begin{tabularx}{\linewidth}{@{}lX@{}}
原函数 $F$ & $F'=f$，要求的是逐点导数关系；\\
定积分 $\int_a^b f$ & 一个数，要求区间上的可积性；\\
变上限函数 $G(x)=\int_a^x f(t)\,dt$ & 先由累积定义函数，再讨论其连续性或可导性。
\end{tabularx}
\end{tcolorbox}
\subsection*{1.\ 两个方向}
若 $f$ 在 $[a,b]$ 连续，定义 $G(x)=\int_a^x f(t)\,dt$，则
\[
G'(x)=f(x),\qquad \int_a^b f(x)\,dx=F(b)-F(a)\quad(F'=f).
\]
第一式是变上限积分的局部结论，第二式是 Newton--Leibniz 的累计结论；不能把一个对象的条件自动套给另一个对象。
\subsection*{2.\ 正则性梯：三种资格不是同一个强弱链}
\[
f\text{ 连续}\Rightarrow G\text{ 可导且 }G'=f,\qquad
f\text{ Riemann 可积}\Rightarrow G\text{ 连续},
\]
但“$f$ 可积”不等于“$G'=f$ 处处成立”，而“$f$ 有原函数”也不等于“$f$ Riemann 可积”。三类对象真正形成的是一个带分叉的资格网络：
\[
\boxed{\begin{array}{c}
f\text{ 连续}\Longrightarrow f\text{ 有原函数}\text{ 且 }f\text{ Riemann 可积},\\[2mm]
f\text{ 有原函数}\not\Longrightarrow f\text{ Riemann 可积},\\
f\text{ Riemann 可积}\not\Longrightarrow f\text{ 有原函数}.
\end{array}}
\]

第二条的生成性反例来自“导数可以无界”。定义
\[
F(x)=\begin{cases}
x^2\sin(1/x^2),&x\ne0,\\
0,&x=0,
\end{cases}
\]
则 $F'(0)=0$，而 $x\ne0$ 时
\[
F'(x)=2x\sin(1/x^2)-\frac{2}{x}\cos(1/x^2),
\]
在 $0$ 附近无界。因此 $f=F'$ 确实有原函数 $F$，却不可能作为闭区间上的 Riemann 可积函数，因为 Riemann 可积首先要求有界。

第三条可用阶跃/符号函数理解：它们有界且只有有限个间断点，因此 Riemann 可积；但若它们存在原函数，就必须作为导函数满足 Darboux 介值性，而跳跃直接违反这一必要条件。

若 $f$ 在 $[a,b]$ Riemann 可积，则它有界，设 $|f|\le M$。于是
\[
|G(y)-G(x)|
=\left|\int_x^y f(t)\,dt\right|
\le M|y-x|.
\]
所以变上限积分不只是连续，而是\textbf{Lipschitz 连续}，从而一致连续。这里的正则化来自“累积把局部振荡先做平均”，并不意味着被积函数本身变得连续。

更精细地，若 $f$ 在 $x_0$ 附近 Riemann 可积，且存在有限极限
\[
\lim_{t\to x_0}f(t)=L,
\]
即使 $f(x_0)\ne L$ 或 $f(x_0)$ 被任意重定义，仍有
\[
G'(x_0)=L,
\qquad G(x)=\int_a^x f(t)\,dt.
\]
原因是单点值不影响积分，而差商
\[
\frac{G(x_0+h)-G(x_0)}{h}
=\frac1h\int_{x_0}^{x_0+h}f(t)\,dt
\]
读取的是邻域平均值；当邻域内 $f(t)$ 统一趋于 $L$ 时，这个平均值也趋于 $L$。因此“变上限函数在点可导”与“$G'(x_0)=f(x_0)$”必须分开判断。

事实上，真正的点态判据就是\textbf{缩小区间上的有向平均值}：
\[
\boxed{
G'(x_0)=A
\iff
\lim_{h\to0}\frac1h\int_{x_0}^{x_0+h}f(t)\,dt=A
}
\]
（前提是这些局部积分有定义）。普通极限 $f(t)\to A$ 只是这个平均值极限成立的一条充分路径，不是必要条件；某些有界快速振荡函数即使点态极限不存在，正负面积仍可能在缩小区间内平均抵消，使 $G'(x_0)$ 存在。

这也解释了为什么“单点修改”对两个世界的影响完全不同。Riemann 积分对有限个点的函数值修改不敏感，但导数关系 $F'(x)=f(x)$ 是逐点条件，改错一个点就可能让 $f$ 不再是导函数。Thomae 函数是一个极端示意：它在有理点 $p/q$（既约）取 $1/q$、无理点取 $0$，Riemann 可积且积分为 $0$；其变上限积分恒为 $0$，所以导数也恒为 $0$，并不等于原来的 Thomae 函数在有理点的值。更一般的“零测集修改不改变 Riemann 积分”属于 Lebesgue 判别的 Extension，本桥核心只需掌握有限点修改即可。

\subsection*{3.\ 原函数存在性的反向约束：导函数不必连续，但必有介值性}

若存在 $F$ 使 $F'=f$，则 $f$ 是导函数。导函数可以不连续，但由 Darboux 定理必具有介值性：若 $x_1<x_2$，则每个介于 $f(x_1)$ 与 $f(x_2)$ 之间的值都由某个中间点取得。于是“有原函数”对间断形态施加强约束：
\begin{itemize}
  \item \textbf{可去间断不可能是导函数的间断。} 若 $\lim_{x\to x_0}f(x)=L$ 存在且 $F'=f$，由导数极限定理（或中值定理）必须有 $f(x_0)=L$；所谓“邻域已经稳定但点值错位”无法发生在真正的导函数上。
  \item \textbf{跳跃间断不可能。} 左右极限若为两个不同有限值，Darboux 介值性要求在任意小邻域内填满两者之间的值，与稳定到两个分离单侧极限矛盾。
  \item \textbf{无穷间断也不可能出现在定义了导数的内部点。} 若例如 $f(x)\to+\infty$（$x\to x_0^+$），则足够靠右时 $f>M$；对 $F$ 在 $[x_0,x_0+h]$ 用 Lagrange，中间点导数大于 $M$，故右差商也大于 $M$，这与有限的 $F'(x_0)=f(x_0)$ 矛盾。
  \item \textbf{振荡间断可能出现，但不是充分条件。} 导函数可以剧烈振荡而不连续；例如
  \[
  F(x)=\begin{cases}x^2\sin(1/x),&x\ne0,\\0,&x=0,\end{cases}
  \]
  在 $0$ 可导，而
  \[
  F'(x)=2x\sin(1/x)-\cos(1/x)\quad(x\ne0),\qquad F'(0)=0
  \]
  在 $0$ 振荡不连续。反过来，一个任意振荡函数若违反介值性，仍不可能成为导函数。
\end{itemize}

因此更可靠的逻辑不是“看间断点名称直接判原函数”，而是
\[
\boxed{\text{若 }f\text{ 有原函数}\Longrightarrow f\text{ 必满足导函数的介值性及点态导数约束}.}
\]
间断分类只是快速暴露哪些形态立即违反这些必要条件。Darboux 性本身仍只是必要条件，不应反向写成“具有介值性就一定有原函数”。

\subsection*{4.\ 单调性：累积函数读的是每个子区间的净面积}
若 $f\ge0$，则
\[
G(y)-G(x)=\int_x^y f(t)\,dt\ge0\qquad(x<y),
\]
故 $G$ 非减。若要升级为\textbf{严格递增}，真正需要的是每个非退化子区间都有正净面积：
\[
\forall x<y,\qquad \int_x^y f(t)\,dt>0.
\]
连续情形下，“$f\ge0$ 且任何子区间都不是恒零”足以保证严格递增；只知道某一个孤立点 $f(c)>0$ 而没有邻域控制，不足以产生正面积。

\subsection*{5.\ Leibniz 公式：上下限移动与被积函数自身变化是三份贡献}
对
\[
H(x)=\int_{\alpha(x)}^{\beta(x)}g(x,t)\,dt,
\]
在 $g,g_x$ 与上下限满足相应连续/可导资格时，
\[
\boxed{
H'(x)
=g(x,\beta(x))\beta'(x)
-g(x,\alpha(x))\alpha'(x)
+\int_{\alpha(x)}^{\beta(x)}g_x(x,t)\,dt
}.
\]
前两项来自\textbf{积分区域边界移动}，最后一项来自\textbf{区域内部每个位置的密度自身变化}。若 $g$ 不显含 $x$，第三项才消失。旧式“上限项减下限项”不能被无条件推广到含参数被积函数。

\subsection*{6.\ 重复累积：$n$ 次积分可以压成一个卷积核}
从同一基点 $a$ 连续做 $n$ 次变上限积分，可压缩成
\[
I_n f(x)
=\frac1{(n-1)!}\int_a^x(x-t)^{n-1}f(t)\,dt.
\]
核 $(x-t)^{n-1}/(n-1)!$ 记录一个局部贡献 $f(t)$ 在后续 $n-1$ 次累积中被覆盖的“组合权重”。在足够正则条件下逐次求导有
\[
(I_n f)^{(n)}=f,
\]
并且 $I_n f$ 在基点自动带有前 $n-1$ 阶零初值。这个接口把“反复积分”与“高阶初值问题”接到 Topic12；一般卷积理论不进入数学一主干。

\subsection*{7.\ 分段原函数的缝合：值能用常数接，导数跳跃接不了}
若 $f$ 在不同连通分支上各自有原函数，可以分别加常数调整函数值；但若接口处 $f$ 的左右行为形成跳跃，任何把两侧原函数拼成连续函数的选择仍无法让其导数在接口处同时等于左右两个不同值。积分常数只移动函数的纵向基线，不能修复导函数的局部结构冲突。

\subsection*{8.\ 反向使用}
遇到 $\int f$：先问当前要的是原函数族还是定积分数值；遇到 $G'(x)$：先把差商翻译成缩小区间平均值；遇到参数积分：区分边界移动项与内部 $g_x$ 项；遇到“有原函数/可积”的互推：先回到 Darboux 与有界性两套不同资格。

\subsection*{9.\ 最小例子与调用}
令 $f(t)=0$（$t<0$），$f(t)=1$（$t\ge0$），并定义 $G(x)=\int_{-1}^{x}f(t)\,dt$。则
\[
G(x)=\begin{cases}0,&x\le0,\\x,&x>0,\end{cases}
\]
$G$ 连续但在 $0$ 不可导，说明“可积 $\Rightarrow G$ 连续”不能升级为“$G'=f$ 处处”。题面混用原函数、定积分和变限函数时调用。
\begin{tcolorbox}[title=检查句]
我当前操作的对象是“一个定积分数值、一个累积函数，还是一个原函数关系”？题目需要的是有界/Riemann 可积、局部平均值极限、Darboux 介值性，还是连续性？如果上下限和被积函数都依赖参数，我是否漏掉了第三个 $\int g_x$ 贡献？
\end{tcolorbox}
\end{document}
